\newpage
\titledquestion{Two Accumulators}

\textbf{Two Accumulators}

\code{foldl} threads one accumulator through a list. Some computations want
\textbf{two}, each able to depend on the other. We will build \code{fold2},
which threads two, and then use it.

\begin{parts}
  \part[8]\hfill

    Implement \code{fold2}:

    \spec
      {fold2}
      {\code{('a * 'b * 'c -> 'b * 'c) -> 'b * 'c -> 'a list -> 'b * 'c}}
      {\code{g} is total}
      {\code{fold2 g (b0, c0) [x1, ..., xn]} threads the two accumulators
      left-to-right: starting from \code{(b0, c0)}, each element \code{xi}
      updates the accumulators by \code{g (xi, b, c)}. It evaluates to the
      final pair of accumulators.}

    That is: like \code{foldl}, except \code{g} receives both accumulators and
    produces both, so either may depend on the other.

    \begin{solutionorbox}[16em]
      \begin{codeblock}
        fun fold2 g (b, c) [] = (b, c)
          | fold2 g (b, c) (x::xs) =
              let
                val (b', c') = g (x, b, c)
              in
                fold2 g (b', c') xs
              end
      \end{codeblock}
    \end{solutionorbox}

  \newpage

  \part[8]\hfill

    The \textit{best prefix sum} is the largest running prefix sum, counting the
    empty prefix as \code{0}. For \code{[3, ~2, ~2, 5]} the prefix sums are
    \code{3, 1, ~1, 4}, so the answer is \code{4}. Implement \code{bestPrefix}:

    \spec
      {bestPrefix}
      {\code{int list -> int}}
      {\code{true}}
      {\code{bestPrefix xs} evaluates to the largest running prefix sum of
      \code{xs}, or \code{0} if every prefix sum is negative.}

    \begin{constraint}
      You must use \code{fold2}, and traverse the list only once. You may not
      use explicit recursion.
    \end{constraint}

    \begin{solutionorbox}[18em]
      The two accumulators are the running \textbf{sum} of the prefix so far,
      and the \textbf{best} prefix sum seen so far. The best cannot be updated
      without the running sum --- that is the interaction.

      \begin{codeblock}
        fun bestPrefix xs =
          #2 (fold2
                (fn (x, sum, best) =>
                   let
                     val sum' = sum + x
                   in
                     (sum', Int.max (best, sum'))
                   end)
                (0, 0)
                xs)
      \end{codeblock}

      We start both accumulators at \code{0} (the empty prefix has sum \code{0},
      and it is a candidate for the best). At each element we advance the sum,
      then take the max of the old best and the \textit{new} sum. The answer is
      the second accumulator.
    \end{solutionorbox}

  \newpage

  \part[4]\hfill

    A student claims \code{fold2} is strictly more powerful than \code{foldl}.
    Settle it by implementing \code{foldl} in terms of \code{fold2}, to show \code{fold2}
    is at least as powerful.

    \begin{solutionorbox}[20em]
      \code{foldl} is the special case of \code{fold2} that ignores the second
      accumulator (we thread \code{unit} through it):

      \begin{codeblock}
        fun foldl' f z xs =
          #1 (fold2 (fn (x, acc, u) => (f (x, acc), u)) (z, ()) xs)
      \end{codeblock}
    \end{solutionorbox}
\end{parts}
